\cotefile{68}
\scanpage{2}
\textcircled{1} Jeu (de position)

\[
\text{\textcircled{a}}\left\{\begin{aligned}
&C\quad\text{ensemble des positions (ou configurations)}\\
&R\subset C\times C\qquad\text{relation}\quad ((x,y)\in R\ \text{se lit }y\in R(x),\ \text{ou}\\
&\hspace{19em}y\ \text{\uline{successeur} de }x\ )
\end{aligned}\right.
\]

\inserted{\uline{Déf} \quad Position $x$ terminale $\Longleftrightarrow R(x)=\varnothing$\\
$C_0=\{x\in C\mid R(x)=\varnothing\}={}$ ens des positions terminales}

\[
\text{\textcircled{b}}\left\{\begin{aligned}
&J\quad\text{ensemble des joueurs}\\
&i:C\setminus C_0\longrightarrow J\qquad i(x)\ \text{est}\ \text{le joueur qui a l'initiative}\\
&\hspace{19em}\text{dans la position }x
\end{aligned}\right.
\]

\margin{\uline{NB} \deleted{$i$ \uline{n'intervient que par sa valeur sur} $C\setminus C_0$ \unsure{$J$}}}
\[
C\setminus C_0=\coprod_{j\in J} C(j)
\]

\[
\begin{aligned}
\text{\textcircled{c}}\quad J&\longrightarrow\mathcal P(C_0)&
j&\longmapsto\underbrace{G_0(j)}_{\substack{\text{positions terminales}\\\text{\uline{gagnantes} \uline{pour} }j}}\subset C_0
\end{aligned}
\]
\margin{\uline{NB} On ne suppose pas $G_0(j)\subset C_0(j)$}

\textcircled{2} \uline{Parties}. \deleted{\unsure{2a}} (Ne dépend que de $(C,R)\ldots$)

\uline{Partie} \inserted{(à $n$ coups, $n\in\mathbb N$)} $(x_0,x_1,\ldots,x_n)\in C^{n+1}$ \qquad $\forall\,0\leq i\leq n-1,\quad x_{i+1}\in R(x_i)$

\uline{Prolongmt} des parties

\uline{Partie \unsure{terminante|terminée}} ($=$ partie maximale) : telles que $x_n\in C_0$

Partie infinie, partie circulaire. \quad $(x_0,x_1,\ldots,x_n)$ \quad où $x_n=x_0$ \quad $(n\geq1)$

\uline{Prop} S'il existe une partie circulaire, existe une partie infinie. L'inverse est vrai, si $C$ est fini.

\[
C_{-1}=\varnothing\qquad C_0=\{x\in C\mid\forall y\in R(x),\,y\in C_{-1}\}=\{x\in C\mid R(x)=\varnothing\}
\]

\uline{Construction} récurrente
\[
\begin{aligned}
C_1&=\{x\in C\mid\forall y\in R(x),\,y\in C_0\} & \text{NB}\quad C_1&\supset C_0\\
&\cdots\\
C_i&=\{x\in C\mid\forall y\in R(x),\,y\in C_{i-1}\}
\end{aligned}
\]

\uline{\unsure{Prop 1|Propt|Prop}} \quad $(C_i)$ \deleted{\unsure{a}} suite croiss. \unsure{de} d'ensembles.
\[
C_i=\{x\in C\mid\text{tte partie \deleted{\unsure{commençant}} \inserted{à partir de} }x\text{ est de longueur }\leq i\}
\]

\uline{Cor} Soit $C_\infty=\bigcup C_i$. Alors $C_\infty=\{x\in C\mid{}$ l'ens des \inserted{longueurs des} parties partant de $x$ est \unsure{borné|bornée}$\}$
% archive sheet number: 1

\scanpage{3}
\uline{Cor} \quad Conditions équivalentes

a) $C=C_\infty$ i.e. pour $\forall x\in C$, l'ens des longueurs des parties partant de $x$ est fini.

b) $\exists\ C\xrightarrow{l}\mathbb N$ tel que si $y\in R(x)$, on ait $l(y)<l(x)$

[ Alors $\forall x\in C$, si $l=l(x)$, on a $x\in C_l$ ]

\deleted{\unsure{Impli.|Exemple}}

\uline{NB} Ces conditions impliquent que $\nexists$ partie infinie. \deleted{Mais} Si $C$ est \deleted{\unsure{in}}fini, \deleted{l'inverse \unsure{est|n'est} \unsure{pas} vraie} \inserted{la réciproque est vraie} : en effet, tte partie \inserted{a tjs. une longueur} \deleted{est} majorée \deleted{\unsure{de}} par $(\operatorname{Card} C)-1=N-1$, i.e. $C=C_{N-1}$

\deleted{\uline{Prop} Soit $C$}

\textcircled{3} \uline{Stratégie pour $j$} $\in J$ \quad $\Sigma\subset(C(j)\times C)\cap R\big|$ tel que
\[
\forall x\in C(j),\quad\text{on a }\quad(R(x)\ne\varnothing\ \text{(i.e. }x\notin C_0)\Longrightarrow\Sigma(x)\ne\varnothing)
\]

Partie $\underline{x}=(x_0,\ldots x_n)$ \uline{compatible} avec stratégie $\Sigma$ pour $j$ :
\[
\forall\,0\leq i\leq n-1\ \text{ tel que }\ x_i\in C(j),\ \text{on a }\ x_{i+1}\in\Sigma(x_i)
\]
\margin{/ Idem pour partie infinie}

\uline{Stratégie} \inserted{(pour $j$)} \uline{gagnante} pour $x$ $(\in C)$ :

a) Toute partie partant de $x$, compatible avec $\Sigma$, est finie \unsure{pour} (i.e. $\nexists$ partie infinie partant de $x$, comp. avec $\Sigma$)

b) Toute partie terminée, \inserted{$(x_0=x,x_1,\ldots,x_{N})$} partant de $x$, compatible avec $\Sigma$, on a $x_N\in G_0(j)$.

On dit que $x\in C$ est \uline{gagnante} \uline{pour} $j$ si $\exists$ stratégie \deleted{gagnante} \inserted{pour} $j$ gagnante pour $x$.

L'ens de ces $x$ est noté $G(j)$

\uline{NB} \qquad $G(j)\cap C_0=G_0(j)$

Posons $\bar C(j)=C\setminus C(j),\quad\bar G(j)=C\setminus G(j)\ldots$
% archive sheet number: 2

\scanpage{4}
\deleted{\uline{Prop}} \inserted{\uline{Thm 1}}
\[
\begin{aligned}
(C(j)\setminus C_0(j))\cap G(j)&=\{x\in C(j)\setminus C_0(j)\mid \deleted{\unsure{x\in G_0(j)}\ \text{ou}}\ \exists y\in R(x),\ \text{avec }\unsure{y\in G(j)}\}\\[-2pt]
&\hspace{11em}\inserted{R(x)\cap G(j)\ne\varnothing}\\
(\bar C(j)\setminus\bar C_0(j))\cap G(j)&=\{x\in\bar C(j)\mid\forall y\in R(x)\ \text{on a }\ y\in G(j)\}\\
&\quad\deleted{\text{et si } R(x)=\varnothing\ \text{i.e. } x\in C_0,\ \text{alors }x\in\unsure{G_0(j)}}\quad\inserted{R(x)\subset G(j)}
\end{aligned}
\]
% slanted
\margin{S'il n'y a pas de partie infinie}

\uline{Dém} \quad \deleted{a) Soit $x\in C(j)$} \inserted{$1^\circ$) \unsure{a}} Soit $x\in C(j)$, il faut prouver que $x\in G(j)\Longleftrightarrow{}$ \deleted{\unsure{l'on a|implique}} $\exists y\in R(x),\,y\in G(j)$.

\deleted{Si $x\in C_0$, c'est clair (cf NB plus haut) \unsure{à vérifier}}

Soit $x\in C\setminus C_0$ \inserted{i.e. $R(x)\ne0$.}, \inserted{Prouvons $\Longrightarrow$ i.e. supposons $x\in G(j)$, prouvons que $\exists y$ tel que\ldots} Soit $\Sigma$ stratégie \deleted{pour $j$} $x$-gagnante. Considérons $\Sigma(x)$ ; donc \deleted{\unsure{Puisque}, on} $\Sigma(x)\ne\varnothing$, donc $\exists y\in\Sigma(x)$. Je dis que $y\in G(j)$, i.e.

\uline{Lemme} \quad Si $x\in C(j)\cap G(j)$, $\Sigma$ une $j$-stratégie $x$-gagnante, alors $\forall y\in\Sigma(x)$, \inserted{on a} $y\in G(j)$. \inserted{i.e. $\Sigma(x)\subset G(j)$.} Plus précisément $\Sigma$ est aussi $y$-gagnante.

\deleted{En effet, je dis que $\Sigma$}

Trivial sur définitions.

Inversement, supposons $\exists y\in R(x)$ avec $y\in G(j)$, prouvons que $x\in G(j)$. Soit $\Sigma$ une $j$-stratégie $y$-gagnante. On va \inserted{essayer de} définir une $j$-stratégie $x$-gagnante \inserted{$\Sigma'$} ainsi : Pour $z\in C(j)$, on pose
\[
\Sigma'(z)=\begin{cases}
\Sigma(z)&\text{si }z\ne x\\
\{y\}\deleted{\unsure{\Sigma(x)}}&\text{si }z=x
\end{cases}
\]

C'est une $j$-stratégie, car \deleted{si \deleted{$R(z)\ne\varnothing$} $R(\unsure{z|x})\ne\varnothing$, alors $\Sigma(z)\ne0$ \deleted{\unsure{pour}} (\unsure{étant donné|on traite} le cas $z\ne x$ pour \unsure{pb}, mais alors $\Sigma'(z)=\Sigma(z)\ne\varnothing$ par hyp. que $\Sigma$ est stratégie.}

Prouvons que $\Sigma'$ est $x$-gagnante : Soit une partie $(x_0=x,x_1,\ldots,x_n\ldots)$ partant de $x$ et compatible avec $\Sigma'$. On a donc $x_1=y$ ; \deleted{et} \inserted{prouvons que} la partie $(\underset{\substack{\Vert\\y}}{x_1},x_2,\ldots x_n\ldots)$ est compatible avec $\Sigma$, i.e. que \deleted{$x_i\in$} $i\geq1,\ x_i\in C(j)\Longrightarrow x_{i+1}\in\Sigma(x_i)$ ;
% archive sheet number: 3

\scanpage{5}
Il n'y a de pb. que si \deleted{\unsure{$\Sigma$}} \inserted{l'un des} $x_i$ \inserted{est égal à} \deleted{$=$} $x$. Car si $x_i\ne x$, $\Sigma(x_i)=\Sigma'(x_i)$, \deleted{\unsure{de|donc}} $x_{i+1}\in\Sigma'(x_i)$ par hyp. de compatibilité \uline{avec $\Sigma'$}.

\deleted{Mais si on avait $x_i=x$, alors $\underset{\substack{\Vert\\y}}{x_1},x_2,\ldots,\underset{\substack{\Vert\\x}}{x_i},\ldots$ on aurait partie circulaire\ldots}

% slanted
\margin{Supposons \inserted{donc d'abord} que \deleted{dans $\Sigma$} il n'y a pas de partie comp. avec \deleted{circulaire} \deleted{\unsure{En a sinon ?|Fermeture ?}} $\Sigma$ commençant avec $y$ et terminant \unsure{avec $x$|par $x$}.}

Donc par hyp. sur $\Sigma$, la partie $(x_1,\ldots x_n\ldots)$ est finie, \deleted{\unsure{de|donc}} et son \deleted{\inserted{dernière} \unsure{position|partie} terminante} \inserted{aboutissement} $x_N$ \inserted{qui} est aussi aboutissement de la partie $(x_0,x_1,\ldots)$, est $\in G_0(j)$, \deleted{cqfd.} OK dans le cas envisagé.

Mais si $\exists$ partie comp. avec $\Sigma$ commençant avec $y$ et terminant avec $x$, alors par le lemme on a \inserted{aussi} $x\in G(j)$ cqfd.

$2^\circ$) Soit $x\in\bar C(j)$ ; prouvons que
\[
\begin{aligned}
x\in G(j)&\Longleftrightarrow\quad\forall y\in R(x),\ \text{on a }\ y\in G(j)\\
&\qquad\deleted{\text{et }(R(x)=\varnothing)\Rightarrow x\in G_0(j)}
\end{aligned}
\]

\deleted{Si $x\in C_0$ i.e. $R(x)=\varnothing$, c'est clair}

\deleted{Si $x\notin C_0$ i.e. $R(x)\ne\varnothing$, il faut prouver $x\in G(j)\Longleftrightarrow\forall y\in R(x),\ y\in G(j)$}

Prouvons $\Longrightarrow$, Soit $x\in G(j)$, $y\in R(x)$, prouvons $y\in G(j)$. Soit $\Sigma$ \inserted{$j$-}stratégie \deleted{\unsure{gag}} \deleted{\unsure{$X$}}$x$-gagnante, je dis qu'elle est $y$-gagnante. En effet considérons \deleted{\unsure{la}} \unsure{$x$-}partie $(\deleted{\unsure{x_0=x}},x_1=y,x_2,\ldots,x_n,\ldots)$ \deleted{\unsure{infinie}}\inserted{maximale}, compatible \deleted{avec $\Sigma$. Mais \uline{NB} comme} \inserted{avec $\Sigma$, alors $(x_0=x,x_1=y,x_2,\ldots x_n,\ldots)$} est aussi comp. avec $\Sigma$ (NB comme $x\in\bar C(j)$, on a $(x,y)$ compatible avec $\Sigma$)

Alors par hyp. sur $\Sigma$, cette partie est finie, (soit $x_N$ son aboutissement) et $x_N\in G_0(j)$. \deleted{\unsure{Pour}} \unsure{Donc} \deleted{prouvons que \unsure{$x_1=y$|$x_1,\ y$} $\in G_0(j)$, on fait une vér. \unsure{déplacement|dépendante} d\unsure{e}, et on est ramené au \uline{lemme} Si $z'\in\Sigma(z)$}
% archive sheet number: 4

\scanpage{6}
% archive sheet number: 5

Prouvons $\Longleftarrow$, i.e. \deleted{Soit $x\in\bar C(j)$, $R(x)\ne\varnothing$,} $\forall y\in R(x)$, on a $\unsure{Y|y}\in G(j)$, prouvons $x\in G(j)$.

Stratégie $\Sigma$ avec :
\[
z\in C(j),\quad \Sigma(z)=
\begin{cases}
 R(z)&\text{si }z\notin G(j)\\
 R(z)\cap G(j)&\text{si }z\in G(j)
\end{cases}
\]

a) \uline{C'est une stratégie}, \deleted{i.e. $R(z)\ne\varnothing\Longrightarrow\Sigma(z)\ne\varnothing$.} \inserted{par 1°. et 2°), $\Rightarrow$}

\deleted{En effet, il n'y a de question que si $z\in G(j)$, et alors $R(z)\cap G(j)\ne\varnothing$ par \unsure{a|$1^\circ$}).}

b) Elle est $x$-gagnante. En effet, soit $(x_0=x,x_1=y,x_2,\ldots x_n,\ldots)$ une partie partant de $x$, comp. avec $\Sigma$, admettons qu'elle est finie.

% slanted
\margin{On suppose qu'il n'y a pas de partie infinie}

Prouvons de proche en proche que $x_i\in G(j)$ $\forall i\geq1$ \deleted{On distingue} \inserted{OK} pour $i=1$ par hyp.

Supposons que \deleted{\unsure{c'est prouvé|c'est p}} l'on ait $x_i\in G(j)$, prouvons $x_{i+1}\in G(j)$. En effet, si $x_i\in\bar C(j)$, cela provient de la partie déjà prouvée de 2°. Si $x_i\in C(j)$, cela provient de la définition de $\Sigma$, car $\Sigma(x_i)=R(x_i)\cap G(j)$ donc $x_{i+1}\in\Sigma(x_i)$ \deleted{\unsure{donc $x_i$|donc $x$|donc $x_{i+1}$}} est $\in\unsure{C|G}(j)$. O.K.

\uline{Corollaire} \deleted{Parmi \unsure{tes|toutes} les stratég} \deleted{Il existe une $j$-stratégie $\Sigma$ telle} \deleted{que} \deleted{Soit $\Gamma\subset G(j)$.} \inserted{Soit $j\in J$, $\Gamma\subset G(j)$ stable (dans $G(j)$) par succession.} a) Il existe une $j$-stratégie $\Sigma$, \unsure{gagnante|gagn.} pour tous les $x\in\Gamma$.

b) Parmi celles-ci il y en a une plus grande

\scanpage{7}
% archive sheet number: 6

\deleted{Prouvons $\Longleftarrow$, i.e. supposons que $\forall y\in R(x)$, on a $y\in G(j)$, prouvons $x\in G(j)$, i.e. montrons qu'il existe une $j$-stratégie $\Sigma$ $x$-gagnante (sachant que \inserted{$R(x)\ne\varnothing$ et que} $\forall y\in R(x)$, $\exists\Sigma_y$ \deleted{\unsure{straté|strati|stratégie}} $j$-stratégie $y$-gagnante, et que $x\in\bar C(j)$).}

\deleted{Considérons un bon ordre sur $R(x)$, et posons \inserted{pour $z\in C(j)$}
\[
\Sigma(\underline{z})=\begin{cases}
R(z)&\text{si }\deleted{\text{$z$ n'est pas conséquent}}\ z\ngtr x\\[-2pt]
 &\inserted{\text{($z$ n'est pas conséquent de $x$)}}\\
\Sigma_{y(z)}(z)&\text{si }\deleted{y(z)}\ z>x,\ \deleted{\text{donc}}\ \text{si }\deleted{z>y}
\end{cases}
\]
$\exists y(z)\in R(x)$, avec $z>y(z)$. On prendra $y(z)$ le \uuline{plus petit}.}

\deleted{On vérifie immédiatement que a) $\Sigma$ \deleted{\unsure{est}} est une \uline{stratégie} i.e. $R(z)\ne\varnothing\Longrightarrow\Sigma(z)\ne\varnothing$ (car les $\Sigma_y$ sont des stratégies), et \inserted{prouvons \deleted{que}} qu'elle est $x$-g. \deleted{\unsure{gagnante|gagnte}}}

\deleted{\deleted{Pour ceci, voyons d'abord que \unsure{si|a} $z\in G(j)\cap C(j)$, \inserted{et si $z$ est un successeur de $x$,} alors $\Sigma(z)\unsure{\subset|\in}G(j)$} \deleted{\unsure{alors}} $\Sigma$}

\deleted{Soit $(x_0=x,x_1=\unsure{y|y_1},x_2,\ldots,x_n,\ldots)$ une \unsure{$\Rightarrow$|$j$} partie partant de $x$, comp. avec $\Sigma$. Alors les $x_i$ $(i\geq1)$ sont \deleted{successeurs de $x$, et on} \inserted{$>x$, donc $\forall i\geq1$, $\exists y_i\in R(x)$} \deleted{voit que si $y_i$ est le p} avec $x_i>y_i$ ; soit $y_i$ le plus petit de ceux ci. Comme $x_{i+1}>x_i>y_i$ donc $x_{i+1}>y_i$, on a $y_{i+1}\leq y_i$, donc les $y_i$ vont en décroissant, donc sont stationnaires. Soit $y$ la valeur commune.}

\deleted{\deleted{Donc} pour $n$ assez grand, on aura}

\scanpage{8}
% archive sheet number: 7

donnée par la condition (pour \deleted{\unsure{$z\notin$|$z\in$|et}} $z\in C(j)$)
\[
\Sigma(z)=\begin{cases}
\text{a)}\ R(z)&\text{si }\deleted{\nexists\text{ partie }(x_0,x_1,\ldots x_n)\text{ avec }x_0\in\Gamma,\ x_n=z}\\
 &\deleted{\nexists x\in\Gamma,\ \unsure{\text{t.q.}\ x_i\in\Sigma(x)|\text{avec }z>x|\text{avec }z\geq x},\ \text{ou si } z\notin G(j)}\\
 &z\notin\Gamma\\
\text{b)}\ R(z)\cap G(j)&\deleted{\text{si }\exists x\in\Gamma\text{ avec }z>x,\ \underline{\text{et}}\ z\in G(j)}\\
 &\deleted{\text{bien (donc }z\in G(j)\text{)}}\\
 &\text{si }z\in\Gamma
\end{cases}
\]

\uline{Dém} Prouvons que $\Sigma(z)$ est une stratégie gagnante pour \unsure{$\forall$|tt|ts} $x\in\Gamma$. C'est une stratégie i.e. $R(z)\ne\varnothing\Longrightarrow\Sigma(z)\ne\varnothing$. Il suffit de regarder le cas \deleted{b), $\exists x\in\Gamma$, avec $z>x$, \uline{et} $z\in G(j)$} \inserted{$z\in\Gamma$} (sinon $\Sigma(z)=R(z)\ne\varnothing$), alors $\Sigma(z)=R(z)\cap G(j)$ et par le 1°) de la prop. on a (puisque $z\notin C_0$ donc $z\notin G_0(j)$) $R(z)\cap G(j)\ne\varnothing$ OK

Prouvons que $\Sigma(z)$ est \uline{gagnante} pour tt $x\in\Gamma$. Soit donc une partie \inserted{$(x_0=x,x_1,\ldots x_N)$} compatible avec $\Sigma$ et maximale, \deleted{\unsure{d'origine|d'originaire|définitive} \unsure{et|en}} \inserted{\unsure{partant|partie} de} $x$, et \unsure{d'aboutissement|destination|d'extrémité|terminée} : $x_N$ ; prouvons $x_N\in G(j)$ (i.e., encore, $x_N\in G_0(j)$ .). Par hyp $x_0\in G(j)$, on va prouver par réc. \inserted{$(0\unsure{\leq\leq|\leq}i\leq N-1)$} $x_i\in G(j)\Rightarrow x_{i+1}\in G(j)$.

\uline{Si} $x_i\in\bar C(\underline{j})$, cela résulte de prop., 2°. Si $x_i\in C(j)$, cela résulte de la définition de $\Sigma$, qui implique que $\Sigma(x_i)=R(z)\cap G(j)$, donc \inserted{comme} $x_{i+1}\in\Sigma(x_i)$, on a $x_{i+1}\in G(j)$. O.K.

Prouvons enfin que parmi les stratégies $\Sigma'$

\scanpage{9}
% archive sheet number: 8

pour $j$ \inserted{$\Gamma$-}gagnantes, \deleted{pour tt} $\Sigma$ est la plus grande, i.e. $\Sigma\supset\Sigma'$, i.e. $\forall z\in C(j)$, on a $\Sigma(\deleted{\unsure{y|x}}z)\supset\Sigma'(\deleted{\unsure{y|x}}z)$. C'est évident (puisque $\Sigma'(z)\subset R(z)$) \deleted{\unsure{si $\Sigma$|si $z$|$\Sigma$}} dans le cas où $\Sigma(z)=R(z)$, donc quand \deleted{on n'a pas simultanément $\exists x\in\Gamma$, tel que $z>x$, \uline{et} $z\in G(j)$.} \inserted{\deleted{\unsure{b)|si}} $z\notin\Gamma$}. On peut supposer \deleted{ces conditions satisfaites} \inserted{$z\in\Gamma$}, donc $\Sigma(z)=R(z)\cap G(j)$. Il suffit donc de prouver que $\Sigma'(z)\subset G(j)$, \deleted{\unsure{i.e.|ici}} ce qui résulte du \uuline{lemme} après la prop.

\deleted{Stratégie}

4°) \uline{Stratégie acceptables}, \uline{positions acceptables}.

$J'$-Stratégie (où $J'\subset J$) \deleted{\unsure{Si|So|Soit}} Posant
\[
C(J')=\bigcup_{j\in J'}C(j),
\]
se définit comme une $j$-stratégie, en y remplaçant $C(j)$ par $C(J')$.

Donc les $j$-stratégies deviennent les $\{j\}$-stratégies.

Une \deleted{\unsure{tt|la}} $J'$-stratégie \inserted{$\Sigma'$} est dite \uline{acceptable} \inserted{pour $x\in C$} \deleted{pour} \uline{contre} (un \uline{$j_0$}$\in J$ \inserted{\unsure{et $x\in C$|distinct}}, si pour \unsure{$\forall$|tt} partie \inserted{$(x_0,\ldots x_N)$} comp. avec $\Sigma$, \deleted{\unsure{aboutissant à|arrivant à}} \inserted{partant de $x$}, et terminée, on a $x_N\notin G(\unsure{j_1|j_0})$ i.e. $x_N\notin G_0(j)$.

\scanpage{10}
% archive sheet number: 9

\deleted{\uline{\unsure{Prop|Prép.|Prop.}} Pour $x\in C$ \unsure{pour|pr|que}}

On dit que \uline{$x$}$\in C$\deleted{\unsure{$(j)$|$(j')$}} est \uline{acceptable} contre $j'\in J$ \deleted{pour $J'=J-\{\unsure{j|j'}\}$} si \deleted{$\exists$} (posant $J'=J-\{j'\}$) $\exists$ $J'$-stratégie $\Sigma$ telle que $\Sigma$ soit acceptable pour $x\in C$ contre $j'$.

% slanted
\margin{Si $x\in C_0$, cela signifie donc que $x\notin G(j')$.}

Cela implique que $x\notin G(j')$. Car si on avait $x\in G(j')$, \deleted{alors \unsure{on|Or}} et si $\Sigma'$ était une $j$-stratégie $x$-gagnante, alors prenant une partie maximale (à priori finie ou infinie) commençant par $x$ et compatible avec $\Sigma$ et avec $\Sigma'$, cette partie serait finie (car $\Sigma$ est \unsure{gagnante|gagn.} pour $x$ et $j$), et on aurait (si $x_N$ est l'aboutissement) à la fois $x_N\in G(j)$ et $x_N\notin G(j)$, absurde. Mieux

\uline{Théorème 2} Soient $j'\in J$, $J'=J\smallsetminus\{j'\}$, $x\in C$.

Pour qu'on ait $x\in\bar G(j')$ (i.e. $x\notin G(j')$), il \uline{faut} et il suffit que $x$ soit acceptable contre $j'$.

\scanpage{11}
% archive sheet number: 10

Il reste à prouver que si $x\in\bar G(j)$, alors $x$ est acceptable contre $j'$. On va construire une stratégie $\Sigma'$ pour $J'$ qui sera acceptable \deleted{contre} \unsure{quant à|contre} $j'$ pour \uline{tous} les $x\in\bar G(j)$.

On posera pour $z\in C(J')$
\[
\Sigma'(z)=\begin{cases}
R(z)&\text{si }z\in G(j')\\
R(z)\cap\bar G(j')&\text{sinon i.e. }z\in\bar C(j')
\end{cases}
\]

\uline{C'est une} $\unsure{J|J'}$ \uline{stratégie} i.e. $R(z)\ne\varnothing\Rightarrow\Sigma'(z)\ne\varnothing$. Il suffit de le voir pour $z\in\bar G(j')$, donc $z\in\bar C(\unsure{j'|J'})\cap\bar G(z')$, mais on a par la prop. 2°
\[
\begin{aligned}
\bar C(j')\cap\bar G(z')=\bigl\{z\in\bar C(j')\bigm|{}&\exists y\in R(z)\text{ avec }y\in\bar G(j')\\
&\text{ou }(R(z)=\varnothing\text{ et }z\in\bar G_0(j'))\bigr\}
\end{aligned}
\]
donc, \unsure{lorsque|comme} $R(z)\ne\varnothing$, on voit que $R(z)\cap\bar G(j')\ne\varnothing$, OK.

Prouvons que $\Sigma'$ est acceptable \deleted{\unsure{p|à}} contre $j'$ pour \unsure{$\forall$|tt} $x\in\bar G(j')$. Soit donc
\[
(\underset{\substack{\rotatebox{90}{$\in$}\\\bar G(j')}}{x_0=x},x_1,\ldots,x_N)\quad\deleted{\text{une}}
\]

\scanpage{12}
% archive sheet number: 11

une partie compatible avec $\Sigma$, commençant avec $x_0\in\bar G(j')$, et maximale i.e. $x_N\in C_0$, prouvons $x_N\in\bar G(j')$. On \deleted{le} prouve par réc. $x_i\in\bar G(j')$\quad $0\leq i\leq N$.

Supposons $x_i\in\bar G(j')$, \inserted{$0\leq i\leq N-1$} prouvons $x_{i+1}\in\bar G(j')\ldots$

À vrai dire, on a fait deux fois le $=$ raisonnement. Considérons un jeu auxiliaire $\mathbb J'=\mathbb J(j')$ : deux joueurs $j'$ et $\overline{j'}=j$, avec $=$ ens des \deleted{situations} positions $C'=C$, \inserted{$R'=R$} et
\[
C'(j')=C(j'),\quad C'(j)=C(J')\quad\inserted{\text{(on posera aussi }C(j)\text{)}}\quad (J'=J-\{j'\}),
\]
\[
\deleted{\text{\unsure{Pour|Puis}}}\quad G'_0(j')=G_0(j'),\quad G'_0(j)=C^{\unsure{\prime|*}}_0\smallsetminus G'_0(j')=C_0\smallsetminus G_0(j')\quad\inserted{\text{(on le note aussi }G_0(j)\text{)}}
\]

\deleted{Alors} $C_0=G'_0(j)\amalg G'_0(j')$. Alors les stratégies \deleted{des} jeu $\mathbb J$ pour $J'=J\smallsetminus\{j'\}$ sont les stratégies du jeu $\mathbb J'=\mathbb J(j')$ pour $j$, et une stratégie \inserted{$\Sigma'$} est acceptable \inserted{pour $x\in C$} contre $j'$ ssi elle est gagnante pour $x\unsure{\in|C}$. Donc les \inserted{positions} $x\in C$ \uline{acceptables} \deleted{pour} \inserted{contre} $j'$ \inserted{(pour $\mathbb J$)} sont celles qui sont gagnantes pour $j$.

% slanted
\margin{s'il n'y a pas de \unsure{partie infinie|parties infinies} \unsure{---|en}}

Comme $G_0(j')\cap G_0(j)=\varnothing$, on sait déjà qu'une $x$ ne peut être simultanément gagnt pour $j$ et $j'$ : \inserted{gagnant pour $j$ $\Rightarrow$ non gagnant pour $j'$} Le raisonnement précédent montre \deleted{que la} l'implication en sens inverse

\scanpage{13}
% archive sheet number: 12

\deleted{Construction}

5. \uline{Essai de construction récurrente \unsure{des|de}} $G(j)$

% slanted
\margin{\uline{NB} $G_{-1}(j)=\varnothing$}

$G_0(j)$ donné
\[
\begin{aligned}
G_1(j)={}&\deleted{\{x\in C}\ G_0(j)\cup\bigl\{x\in C\bigm|{}&
 (x\in C(j)\text{ et }\deleted{\unsure{\exists y|y}}\, R(x)\cap G_0(j)\ne\varnothing)\\
&&\inserted{\text{ou}}\\
&&x\in(\bar C(j)\smallsetminus\bar C_0(j))\text{ et } R(x)\subset G_0(j)\bigr\}\\
\ldots\\
G_{\unsure{i|2}}(j)={}&G_0(j)\cup\bigl\{x\in C\bigm|{}&
 x\in C(j)\text{ et }R(x)\cap G_{i-1}(j)\ne\varnothing\\
&&x\in(\bar C(j)\smallsetminus\bar C_0(j))\text{ et }R(x)\subset G_{i-1}(j)\bigr\}
\end{aligned}
\]

On voit par récurrence sur $i$
\[
G_i(j)\subset G_{i+1}(j)
\]
et aussi
\[
G_i(j)\subset G(j)
\]
(en utilisant la prop.)

Posons
\[
G_\infty(j)=\bigcup G_i(j)\qquad\text{donc}\qquad G_\infty(j)\subset G(j)
\]

? \uline{Th 3} $G_\infty(j)=G(j)$.

% slanted
\margin{Si \uline{$C$ fini} et si \unsure{n'existe de|n'existe|n'existe pas de} partie infinie. Plus généralemt si $C=\bigcup C_i$ cf § 2}

\deleted{Il reste à prouver $x\in G(j)\Longrightarrow\exists i$, $x\in G_i(j)$ :}

\deleted{Soit $\Sigma$ une $j$-stratégie gagnante pour $x$. Donc $\exists$ partie \inserted{finie} maximale \deleted{commençant} \inserted{partant} de $x$, comp. avec $\Sigma$
\[
X=(x_0=x,x_1,\ldots,x_N)
\]
et $x_N\in G_0(j)$. Je dis que \unsure{Voir la p|Vue la f|Vu la f|On le p}}

On prouve par récurrence sur $i\geq0$
\[
G_\infty(j)\cap C_i=G(j)\cap C_i
\]
en utilisant le th. 1.

\scanpage{14}
% archive sheet number: 13

6. \uline{Elimination des parties \deleted{\unsure{circulatoires|circulaires}} \inserted{infinies}} (cas $C$ fini)

\scanpage{15}
% archive sheet number: 14

\[
\begin{aligned}
\{0\}\ \big/\ &\overset{\deleted{\unsure{x|8}}}{\underset{\substack{\rotatebox{90}{$\in$}\\x_1}}{V_1}}\ \big/\ 
\overset{V}{\deleted{x}V_1+\{x_2\}},\ \underset{\substack{\rotatebox{90}{$\in$}\\x_2}}{V_2}\ \big/\ V_2+\{x_3\}\ \deleted{\unsure{\big/|\}|\{}},\\
&V_2+x_3\ast{}V_1,\ V_2+x_3\text{ et }V_1+\{x_2+x_3\},\\
&\hfill V_3
\end{aligned}
\]

\[
\begin{aligned}
\deleted{\unsure{g\quad a\quad b|9\quad a\quad1|q\quad a\quad t}}&\\
 a_1\ a_2\ a_3\quad b_1\ b_2\ b_3&\qquad V_1+\{x_2\}\qquad V_2+\{\deleted{y}x_3\}\\
&3\,V_2+x_3\ast V_1+\{\deleted{\unsure{V|x}}x_2+x_3\}\\
&3\ 3\ 3\ 3\ 1\ 1\ 1
\end{aligned}
\]

\[
E\mathbin{\dot+}x\text{ et }E
\]

\[
(E_{2n}\ \unsure{\text{et}|+|\supsetneq}\ E_2)\ \text{\unsure{tressés|itérés|donnés} ssi }\left[\begin{aligned}
&E_2\ \text{ss-gpe (\unsure{élémtaire 2|élémtre ?|abélien ?})}\\
&E_{2n}\ \text{stable par translation de }E_2
\end{aligned}\right.
\]

\deleted{$n$ pair,}

\[
E_{2n+1}, E_2\ \text{\unsure{tressé|torsé|donnés}}\qquad
\left\{\begin{aligned}
&E_2\ \text{ss-gpe}\\
&E_{2n+1}=E'_{2n}+\{x\},\text{ où }E'_{2n}\text{ est}\\
&\hspace{4em}\text{stable par }E_2
\end{aligned}\right.
\]

\[
E_3\quad E_3\quad\text{tressé}\qquad3\ 2\ 2\ 2\qquad 3\ \unsure{2\ 2|2\ 7}\ 1\ 1
\]

\[
\begin{aligned}
a_1\ a_2\ a_3\quad\unsure{\cdot|\prime\prime}\quad b_1\ b_2\ b_3&\qquad\ldots\\
0\ a\ a'\qquad0\ a\ b'&
\end{aligned}
\]

\diagram{\(\begin{array}{ccc@{\qquad}l@{\qquad}lll}
0\tikz[overlay]{\draw (0,.18) .. controls (1.9,1.2) and (.65,-.4) .. (-.2,-.48) .. controls (-.4,-.42) and (-.03,-.12) .. (.47,.23);} &a&\deleted{b'}&\deleted{\unsure{\notin|\ne}} &0\ a\ a'&0\ a&\underline{a-a'}\\
a&\deleted{2a}&\deleted{a+b'}&\deleted{2a\ne a}&0\ a\ 2a&0\ a&-a\\
\deleted{a'}&\deleted{a+a'}&a'+b'&\deleted{\unsure{a'+b'|a'b'}=a'=b'=a}\\
&&&\deleted{a'+b'}\\
&&\deleted{a'+b'=a}&\deleted{b'=a}\\
&&a'\unsure{+|-}b'=a&&0\\
&&&&a
\end{array}\)}

\[
\underset{\tikz[baseline]{\draw (0,.12)--(0,0)--(2.25,0)--(2.25,.12);}}{-2a\quad-a\quad0\quad a}\quad2a
\qquad \deleted{\unsure{a\ne0|a\ne a}}\quad 2a\ne0
\]

\diagram{\(\begin{array}{ccc}
0\tikz[overlay]{\draw (-.30,-.18) .. controls (.16,-.04) and (.90,.64) .. (1.04,.12) .. controls (1.10,-.32) and (1.23,-.55) .. (1.67,-.54) .. controls (2.67,-.60) and (2.47,-1.22) .. (1.90,-.99) .. controls (1.40,-.98) and (1.32,-.13) .. (.84,-.18) .. controls (.50,-.20) and (-.10,-.82) .. (-.28,-.67) -- (-.34,-.31);} & a &\underline{a-a'}\\
a&\underline{2a}&2a-a'\\
\underline{a'}&\underline{a+a'}&a
\end{array}\)}

\[
\begin{aligned}
&\deleted{\unsure{2a=a|2a-a|2a=0}}\qquad\deleted{2a=0}\\
0\quad\big|&\quad 2a=0\qquad\deleted{\unsure{a-a'=0|a+a'=0}}\\
\big|&\quad a'=-a\\
\big|&\quad a'=2a
\end{aligned}
\]

\diagram{\(\begin{array}{ccc@{\quad}ccc}
\deleted{0}&\deleted{a}&\deleted{-a}&\deleted{0}&\deleted{a}&\deleted{2a}\\
-a&\deleted{0}&\deleted{a}\\
-a&0&a
\end{array}\qquad\begin{bmatrix}
-2a&-a&0\\-a&0&a\\0&a&2a
\end{bmatrix}
\tikz[overlay]{\draw[->] (-.45,-.68) .. controls (-1.0,-1.58) and (-1.45,-1.84) .. (-2.38,-.86) -- (-1.82,-.18);}
\)}

\scanpage{16}
% archive sheet number: 15

\deleted{``cellules'' des décompositions cellulaires ``à la Bruhat'' des espaces homogènes projectifs sous des groupes semi-simples. Le cas particulier des ``cycles de Schubert'' sur les variétés grassmaniennes permet justement de développer des conséquences d'intérêt général en topologie différentielle.}

\[
\begin{aligned}
&&x^{\unsure{2|\prime}}=u\\
(V_1\mathbin{\dot+}\{x\})\ast(V_1\mathbin{\dot+}\{x\})&=2V_1\mathbin{\dot+}\{u\}+2\{x\}\ast V_1\\
&=2G+\{u\}
\end{aligned}
\]

\[
\{e\}\quad\underset{\tikz[baseline]{\draw (0,.13)--(0,0)--(1.65,0)--(1.65,.13);}}{V_1\ \overset{\inserted{V_1\mathbin{\dot+}\{x\}}}{\deleted{\unsure{\{x\}+V_1|\{x\}\ V_1}}}}\quad V_2\qquad
\deleted{2V_1+\unsure{u|x}+}\qquad\boxed{3\ 2\ 2\ 2}
\]

\[
\begin{aligned}
&&E\unsure{\text{et}|\ast}F\qquad a+b&=a'+b'\\
a_1\mathbin{\dot+}\deleted{b}a_2\qquad b_1\mathbin{\dot+}b_2&& (a-a')&=(b'-b)\\
\deleted{a_1\ \unsure{\cdot|i}}\qquad a&&0\quad a\quad\big/\quad0\quad\deleted{\unsure{b+a|b\quad a}}a\\
-\qquad&&\deleted{0\quad a}\quad0\quad a,a\quad2a\qquad a\ne0\\
&&&\deleted{2a\ne\unsure{a|0}}\\
&&&2a=0
\end{aligned}
\]

\diagram{\(\begin{array}{llll@{\qquad}llll}
1&1&1&1&2&2\\
2&1&1&&2&1&1\\
&&&&\deleted{2}\unsure{1|4}&1&1&1\\[3pt]
&&&&2&2&1&1\\
&&&&\multicolumn{4}{l}{1\quad1\quad1\quad1\quad1\quad1}\\
&&&&2&1&1&1\\
&&&&\multicolumn{4}{l}{\deleted{\unsure{2\quad1\quad1|2\quad2\quad1|2\quad2\quad2}}}
\end{array}\)}

\[
\begin{aligned}
\big/\qquad V_1\ast(V_1+\{x\})&=2V_1+x\ast V_1\\
&\underline{V_1+G}\\
E_3+x\ast E_3&\\
&\deleted{\unsure{E+x\ast E|E+x+E}}
\end{aligned}
\]

